% Generated by scripts/assemble.py; corrections belong in source/corrections.json.
\examhead{46}{上海大学 2025 年数学分析试卷}{48}
% p049-b05; source PDF page 49
\par\Needspace{5\baselineskip}
\blockanchor{p049-b05}
\majorquestion{一}计算题. 每题 10 分, 共 40 分.\par
% p049-b06; source PDF page 49
\blockanchor{p049-b06}
\question{1}求数列极限: (1) $\lim_{n\to \infty}\sqrt[n]{n}$, (2) $\lim_{n\to \infty}(2n!)^{\frac{1}{n^2}}$.\par
% p049-b07; source PDF page 49
\blockanchor{p049-b07}
\question{2}设函数 $f(x)$ 在 $[-1, 1]$ 上连续，求 $\lim_{x \to 0} \frac{\ln(1 + f(x) \arctan x)}{2^x - 1}$.\par
% p049-b08; source PDF page 49
\par\Needspace{6\baselineskip}
\blockanchor{p049-b08}
\question{3}求曲线积分\par
\nopagebreak[4]
% p049-b09; source PDF page 49
\blockanchor{p049-b09}
\displayformula{\oint_ {L} \frac {x ^ {2} y}{(x - a) ^ {2} + y ^ {2}} \mathrm{d} y.}
% p049-b10; source PDF page 49
\blockanchor{p049-b10}
其中 L 是由圆  $(x-a)^{2}+y^{2}=a^{2}(a>0)$  与 x 轴围成的第四象限区域的边界, 方向为逆时针.\par
% p049-b11; source PDF page 49
\par\Needspace{6\baselineskip}
\blockanchor{p049-b11}
\question{4}计算曲面积分\par
\nopagebreak[4]
% p049-b12; source PDF page 49
\blockanchor{p049-b12}
\displayformula{\iint_ {\Sigma} x ^ {2} \mathrm{d} y \mathrm{d} z + y \mathrm{d} z \mathrm{d} x + z \mathrm{d} x \mathrm{d} y.}
% p049-b13; source PDF page 49
\blockanchor{p049-b13}
其中 $\Sigma$ 为曲面 $z = \sqrt{x^2 + y^2}$ 在 $0 \leq z \leq 1$ 的部分, 方向取下侧.\par
% p049-b14; source PDF page 49
\par\Needspace{5\baselineskip}
\blockanchor{p049-b14}
\majorquestion{二}证明题. 前 6 道题每题 15 分, 最后 1 道 20 分, 共 110 分.\par
% p049-b15; source PDF page 49
\blockanchor{p049-b15}
\question{1}若函数 $f(x)$ 为区间 $[a, b]$ 上只有有限个间断点的有界函数, 证明: $f(x)$ 在 $[a, b]$ 上可积.\par
% p049-b16; source PDF page 49
\par\Needspace{6\baselineskip}
\blockanchor{p049-b16}
\question{2}设函数 $f(x)$ 在 $[0,1]$ 上二阶可微, $f(0) = f(1) = 0$, $\min_{0 \leq x \leq 1} f(x) = -1$. 证明:\par
\nopagebreak[4]
% p049-b17; source PDF page 49
\blockanchor{p049-b17}
\displayformula{\min _ {0 \leq x \leq 1} f ^ {\prime \prime} (x) \leq 8, \max _ {0 \leq x \leq 1} f ^ {\prime \prime} (x) \geq 8.}
% p049-b18; source PDF page 49
\blockanchor{p049-b18}
\question{3}对于级数 $\sum_{n=1}^{\infty} (-1)^{n} \frac{x^{2} + 2^{n}}{(n+1)!}$.\par
% p049-b19; source PDF page 49
\blockanchor{p049-b19}
(1) 证明: 在任意有限区间 I 上一致收敛.\par
% p049-b20; source PDF page 49
\blockanchor{p049-b20}
(2) 判断级数在任何一点的敛散性, 若收敛, 说明是条件收敛还是绝对收敛, 并给出理由.\par
% p049-b21; source PDF page 49
\blockanchor{p049-b21}
\question{4}解答如下问题:\par
% p049-b22; source PDF page 49
\blockanchor{p049-b22}
(1) 叙述  $R^{2}$  的聚点定理.\par
% p049-b23; source PDF page 49
\blockanchor{p049-b23}
(2) 利用聚点定理证明二元函数 $f$ 在有界闭域 $D \subset \mathbb{R}^2$ 连续则一致连续.\par
% p049-b24; source PDF page 49
\blockanchor{p049-b24}
\question{5}设函数 $\varphi(x)$ 在 $[a, b]$ 上有界可积, 证明: $\lim_{p \to +\infty} \int_{a}^{b} \varphi(x) \sin px \, \mathrm{d}x = 0$.\par
% p049-b25; source PDF page 49
\blockanchor{p049-b25}
\question{6}设函数 $f(x, y) = \begin{cases} \frac{xy^2}{\sqrt{x^2 + y^2}}, & (x, y) \neq (0, 0); \\ 0, & (x, y) = (0, 0). \end{cases}$\par
% p049-b26; source PDF page 49
\blockanchor{p049-b26}
(1) 计算  $f(x,y)$  的偏导数.\par
% p049-b27; source PDF page 49
\blockanchor{p049-b27}
(2) 讨论 $f(x, y)$ 在原点的连续性.\par
% p049-b28; source PDF page 49
\blockanchor{p049-b28}
(3) 讨论 $f(x, y)$ 在原点的可微性.\par
% p050-b00; source PDF page 50
\blockanchor{p050-b00}
\question{7}解答如下问题:\par
% p050-b01; source PDF page 50
\blockanchor{p050-b01}
(1) 若函数 $f(x)$ 满足以下条件: (i) 函数 $f(x)$ 在点 $x_0$ 的邻域 $U(x_0)$ 连续; (ii) 在 $U^\circ(x_0)$ 可导; (iii) 极限 $\lim_{x \to x_0} f'(x)$ 存在. 证明: $f(x)$ 在点 $x_0$ 处可导, 且 $f'(x_0) = \lim_{x \to x_0} f'(x)$.\par
% p050-b02; source PDF page 50
\blockanchor{p050-b02}
(2) 若函数 $f(x)$ 只满足本题 (1) 中条件 (i) 和 (ii), 但不满足 (iii), 则 $f(x)$ 在点 $x_0$ 不可导, 是否正确? 请给出理由.\par
